\documentclass[14pt,a4paper,oneside]{extbook}
\usepackage[utf8]{vietnam}
\usepackage{amsmath,amssymb,mathrsfs,fancyhdr,enumerate,multirow,makecell}
\usepackage{tkz-euclide,tikz-3dplot,tikz,tkz-tab}
\usetikzlibrary{arrows.meta,calc,intersections,angles,quotes,backgrounds,shapes.geometric,patterns,through,shadows,positioning}
\usepackage{yhmath}
\usepackage[top=0.7cm, bottom=0.7cm, left=1.0cm, right=1.0cm, headheight=16pt, headsep=3mm, footskip=5mm]{geometry}
\usepackage[hidelinks,unicode]{hyperref}
\usepackage{qrcode}
\usepackage[many]{tcolorbox}
\usepackage{fontawesome}

\pagestyle{fancy}
\fancyhf{}
\fancyhead[L]{\footnotesize\bfseries\color{teal!85!black}\faBook~SỔ TAY KIẾN THỨC TOÁN 4}
\fancyhead[R]{\footnotesize\bfseries\color{blue!85!black}Lớp Toán Thầy Lĩnh -- loptoanthaylinh.vn}
\fancyfoot[C]{\footnotesize\bfseries\color{gray!80!black}-- Sổ tay kiến thức: Bài 05. Giải bài toán có ba bước tính --}
\renewcommand{\headrulewidth}{0.4pt}

\begin{document}

% --- BANNER TIÊU ĐỀ BÀI HỌC ---
\noindent
\begin{tcolorbox}[
  enhanced,
  colback=teal!8!white,
  colframe=teal!80!black,
  arc=3mm, boxrule=1pt,
  left=3mm, right=3mm, top=2mm, bottom=2mm,
  halign=center
]
  {\Large\bfseries\color{teal!90!black}\faBookmark\ BÀI 05: GIẢI BÀI TOÁN CÓ BA BƯỚC TÍNH}\\[1.5pt]
  {\footnotesize\itshape\color{black!75} Quy trình phân tích đề, vẽ sơ đồ đoạn thẳng và tìm các đại lượng trung gian}
\end{tcolorbox}

\vspace{2.5mm}

% --- KHỐI 1 & 2: LÝ THUYẾT & KĨ NĂNG ---
\noindent
\begin{minipage}[t]{0.49\textwidth}
  \begin{tcolorbox}[
    enhanced,
    title={\footnotesize\bfseries\color{white}\faSearch\ 1. Bản chất bài toán 3 bước},
    colback=blue!2!white,
    colframe=blue!75!black,
    arc=2.5mm, boxrule=0.8pt, top=1.8mm, bottom=1.8mm, left=2mm, right=2mm
  ]
    \footnotesize
    $\bullet$ \textbf{Đặc điểm:} Để trả lời câu hỏi chính của bài toán, ta không thể tính ngay trong 1 phép tính mà phải qua \textbf{2 phép tính trung gian}.\\
    $\bullet$ \textbf{Các mô hình hay gặp:}\\
    - Tìm đại lượng thứ hai $\rightarrow$ tìm đại lượng thứ ba $\rightarrow$ tìm tổng cả ba.\\
    - Tính số lượng $\rightarrow$ tính số tiền $\rightarrow$ tính số tiền còn lại.
  \end{tcolorbox}
\end{minipage}%
\hfill
\begin{minipage}[t]{0.49\textwidth}
  \begin{tcolorbox}[
    enhanced,
    title={\footnotesize\bfseries\color{white}\faCheckSquareO\ 2. Kĩ năng giải chuẩn mực},
    colback=teal!2!white,
    colframe=teal!75!black,
    arc=2.5mm, boxrule=0.8pt, top=1.8mm, bottom=1.8mm, left=2mm, right=2mm
  ]
    \footnotesize
    $\bullet$ \textbf{Lời giải rõ ràng:} Mỗi bước tính phải có lời giải phù hợp với câu hỏi.\\
    $\bullet$ \textbf{Đơn vị đo (danh số):} Phải ghi trong dấu ngoặc đơn $( )$.\\
    $\bullet$ \textbf{Viết gộp biểu thức:} Khi làm quen, có thể viết gộp 3 phép tính thành 1 biểu thức.\\
    $\bullet$ Luôn có dòng \textbf{Đáp số:} ở cuối bài.
  \end{tcolorbox}
\end{minipage}

\vspace{2.5mm}

% --- KHỐI 3: SƠ ĐỒ TƯ DUY ---
\noindent
\begin{tcolorbox}[
  enhanced,
  title={\footnotesize\bfseries\color{white}\faSitemap\ Sơ đồ tư duy: Tiến trình 3 bước giải bài toán có lời văn},
  colback=yellow!2!white,
  colframe=teal!85!black,
  arc=3mm, boxrule=0.9pt, top=1.8mm, bottom=1.8mm, left=1.5mm, right=1.5mm,
  halign=center
]
  \begin{tikzpicture}[
    scale=0.86,
    every node/.style={transform shape},
    root/.style={rectangle, rounded corners=5pt, fill=teal!85!black, text=white, font=\bfseries\footnotesize, inner sep=4.5pt, align=center},
    stepBox/.style={rectangle, rounded corners=4pt, fill=blue!10, draw=blue!75!black, text=blue!90!black, font=\scriptsize, inner sep=4pt, align=center}
  ]
    \node[root] (core) at (0, 0.85) {\faQuestionCircle\ TIẾN TRÌNH 3 BƯỚC GIẢI TOÁN CÓ LỜI VĂN};

    \node[stepBox] (n1) at (-5.0, -0.15) {\textbf{BƯỚC TÍNH 1}\\[1.5pt]Tìm đại lượng trung gian 1};
    \node[stepBox] (n2) at (0, -0.15) {\textbf{BƯỚC TÍNH 2}\\[1.5pt]Tìm đại lượng trung gian 2};
    \node[stepBox] (n3) at (5.0, -0.15) {\textbf{BƯỚC TÍNH 3}\\[1.5pt]Trả lời câu hỏi chính (Đáp số)};

    \draw[teal!85!black, thick, -{Stealth[length=2mm]}] (n1.east) -- (n2.west) node[midway, above, font=\tiny\bfseries] {tiếp theo};
    \draw[teal!85!black, thick, -{Stealth[length=2mm]}] (n2.east) -- (n3.west) node[midway, above, font=\tiny\bfseries] {tiếp theo};
  \end{tikzpicture}
\end{tcolorbox}

\vspace{2.5mm}

% --- KHỐI 4: MẸO & BẪY ---
\noindent
\begin{minipage}[t]{0.54\textwidth}
  \begin{tcolorbox}[
    enhanced,
    title={\footnotesize\bfseries\color{white}\faKey\ Kĩ năng tóm tắt bằng sơ đồ đoạn thẳng},
    colback=green!3!white,
    colframe=green!60!black,
    arc=2.5mm, boxrule=0.8pt, top=1.8mm, bottom=1.8mm, left=2.2mm, right=2.2mm
  ]
    \footnotesize
    $\bullet$ Đọc kĩ các từ khóa: \textit{"nhiều hơn", "ít hơn", "gấp số lần", "giảm đi số lần"}.\\
    $\bullet$ Vẽ các đoạn thẳng có vạch xuất phát thẳng hàng nhau để dễ dàng quan sát sự chênh lệch.\\
    $\bullet$ Điền đầy đủ dấu hỏi chấm $(?)$ vào các đại lượng cần tìm.
  \end{tcolorbox}
\end{minipage}%
\hfill
\begin{minipage}[t]{0.44\textwidth}
  \begin{tcolorbox}[
    enhanced,
    title={\footnotesize\bfseries\color{white}\faExclamationTriangle\ Mẹo tránh bẫy trắc nghiệm},
    colback=red!2!white,
    colframe=red!70!black,
    arc=2.5mm, boxrule=0.8pt, top=1.8mm, bottom=1.8mm, left=2.2mm, right=2.2mm
  ]
    \footnotesize
    $\bullet$ \textbf{Bẫy dừng lại ở bước 2:}\\
    Nhiều học sinh tìm xong đại lượng trung gian tưởng là xong bài, quên mất câu hỏi chính của đề!\\
    $\bullet$ \textbf{Bẫy "gấp" vs "nhiều hơn":}\\
    - Gấp $n$ lần: làm phép nhân ($\times$).\\
    - Nhiều hơn $n$ đơn vị: làm phép cộng ($+$).
  \end{tcolorbox}
\end{minipage}

\end{document}
